\section{Conclusions}
\label{sec:conclusions}

A controlled ablation was conducted to measure segmentation-guided synthetic augmentation for rotated vehicle detection in Peruvian traffic footage. The process included an audit of SMART annotations, conversion to YOLO-OBB labels, preservation of a clip-separated validation set, generation of 1,501 training-only image--label pairs with 3,391 inserted objects, training of YOLO26s-OBB configurations, and uniform evaluation of six artifacts through Macro AP-rIoU and homography-based motion filtering.

Base 1 provides the best aggregate result, with 0.6917 Macro AP-rIoU and 0.7662 AP@0.50. Base 2 is close at 0.6892 and leads at the strictest reported overlap level, with AP@0.80 equal to 0.5417. Improvement C is the strongest proposed alternative, obtaining 0.6682 Macro AP-rIoU while also reaching the highest throughput of 98.93 FPS. Its advantage over Improvements A and B does not extend to either raw-image baseline. The tested SAM-assisted copy-paste strategies therefore offer no accuracy gain over the best conventional training condition in this experiment.

At category level, Motocicleta remains the principal weakness, since every trained run records mean AP-rIoU below 0.30 for that label. The comparison between Base 1 and Base 2 also reveals different error profiles: Base 1 suppresses false detections more effectively and leads the aggregate measure, whereas Base 2 recovers more objects and performs better under strict overlap. These differences reinforce the need to judge augmentation with class-aware rotated metrics and explicit error counts, rather than with generated-sample volume or training-time mAP alone.

The main contribution is an auditable protocol for testing synthetic OBB data without contaminating validation. It shows that the combined proposed configuration may be attractive when processing speed is prioritized, but it is not the most accurate detector. Stronger evidence will require repeated seeds, augmentation targeted at difficult small classes, systematic inspection of detection errors, and comparison against model-level improvements designed specifically for oriented boxes.
